\documentclass[12pt]{article}
\usepackage{amsmath,amssymb}
\usepackage{geometry}
\usepackage{hyperref}
\usepackage{setspace}
\geometry{margin=1in}
\setstretch{1.2}

\title{A Wave Curvature Equation for Unification\\
\large Part One: The Gravitational Potential Pressure Wave}
\author{C.R. Kunferman}
\date{}

\begin{document}

\maketitle

\section*{Abstract}

A mathematical comparison of measurements and calculations.

The relationship between various fields can first be seen when drawing up a comparison chart of shared properties. There exist not only similar correlations, but in some cases identical forms of computation used to quantify and study some of the basic fundamentals.

These comparisons do not imply that the forces are related in terms of occurrence or functionality, they are simply to illustrate the ways we can observe them in similar ways, to find commonalities among them. The comparisons will allow us to discern key differences, and also key areas of likeness in which we could apply similar methods and experimentation towards.

To begin our comparison we take a look at two well known and observable phenomena with an abundance of similar attributes.

\section*{Key Commonalities in the Comparison of Attributes (Gravity, Sound)}


\begin{table}
    \centering
    \begin{tabular}{ccc}
        \textbf{Attribute} & \textbf{Gravity} & \textbf{Sound} \\
Propagation & Gravitational Waves & Longitudinal Waves \\
Speed & Speed of Light & Speed of Sound \\
Frequency & Gravitational Wave Frequency & Acoustic Frequency \\
Amplitude & Strength of gravitational field & Sound pressure level \\
Attenuation & Inverse Square Law & Weakens with distance \\
Detection & Laser Interferometer & Laser Doppler Vibratory \\
Source & Massive object collisions & Vibrating objects \\
Wave Type & Transverse & Longitudinal \\
Energy Transfer & Transfers gravitational energy & Transfers acoustic energy \\
Mathematical Description & General Relativity & Hooke's Law \\
    \end{tabular}
    \caption{Gravity and Sound Comparison}
    \label{tab:Comparison Table 1}
\end{table}


In the above comparison chart, we find that sound and gravity share some key components of behavior: Wave Propagation, Frequency, Inverse Square Law, and Energy transfer. Both can be measured and detected with the use of lasers as well.

Given the abundance of correlations, we then look to the calculations and to that of the mathematical descriptions for both, and here we find correlations aside from the systems of measurements. By drawing a detailed comparison chart containing calculations it presents additional correlations on calculative level to which we apply our measurements and make out predictions with great accuracy.

\section*{Calculations of Gravity vs. Sound}

\begin{align*}
\text{General Calculations:} \quad & G_{\mu\nu} + \Lambda g_{\mu\nu} = \frac{8 \pi G}{c^4} T_{\mu\nu} \\
& \Delta P = \Delta P_{\max} \sin(kx \mp \omega t + \varphi)
\end{align*}

\begin{align*}
\text{Frequency:} \quad & F_g = \frac{\omega}{2\pi} \\
& F = \frac{v}{\lambda}
\end{align*}

\begin{align*}
\text{Energy:} \quad & \nabla P E = mgh \\
& P = E \lambda T
\end{align*}

\begin{align*}
\text{Amplitude:} \quad & h \propto \frac{1}{r} \\
& A \propto \frac{1}{r}
\end{align*}

\begin{align*}
\text{Speed of Propagation:} \quad & c \approx 3 \times 10^{8} \text{ m/s} \\
& c \approx 343 \text{ m/s}
\end{align*}

\begin{align*}
\text{Detection:} \quad & \frac{\delta L}{L} \\
& \delta P \quad \text{(Laser Detection)}
\end{align*}

\begin{align*}
\text{Waveform:} \quad & h(t) \\
& p(t)
\end{align*}

\begin{align*}
\text{Attenuation:} \quad & \frac{1}{r^{2}} \\
& \frac{1}{r^{2}}
\end{align*}

Of particular interest: Many of the equations are so similar that variables can be swapped, or cross applied between these fundamentally different fields, but still hold the same principle calculations.

For instance, by taking Hooke's law and applying it to gravity, we can determine the position of the first object given the starting position of where the gravity is emitted. Likewise, we can apply Newton’s Law of gravitation to sound force, and thanks to the Inverse Square Law they share we are able to successfully calculate how much force it takes for a sound wave to knock down a wall. The base equations would remain the same, while the interpretation and exact amounts calculated would differ.

\subsection*{Using Newton’s Gravitational Law for Sound:}

\begin{align*}
F &= \text{Force based on an inverse square law with respect to distance} \\
G &= \text{A proportionality constant} \\
m_1 &= \text{Density of first object} \\
m_2 &= \text{Density of second object} \\
r &= \text{Distance between objects}
\end{align*}

\begin{align*}
\text{Gravity:} \quad & F = G \frac{m_1 m_2}{r^2} \\
\text{Sound:} \quad & F = C \frac{p_1 p_2}{r^2}
\end{align*}

Likewise, we can do a similar substitution using Hooke's law for calculating the force of gravity:

\begin{align*}
\text{Sound:} \quad & F = -k x \\
\text{Gravity:} \quad & F = -G x
\end{align*}

The dimensional correlations continue with Amplitude, as does the similarity in Waveform, again with only the variables, scale and our interpretations providing the difference. Fundamentally and mathematically however, the numbers crunch the same. Here we will show the two separated by parentheses and the calculation variables or parts they share in brackets.

\begin{align*}
\text{Amplitude:} \quad & (A)(h) \left[\frac{1}{r}\right] \\
\text{Wave Form:} \quad & (p)(h)[t]
\end{align*}

With Wave Equation, Frequency, Amplitude, Waveform and Attenuation all showing correlation, it hints towards a common underlying fundamental force at work, and suggests the two forces may be of the same characteristic composition.

The similarities in the behavior of both occurring on two completely different scales and speeds are striking enough to tempt the observer into thinking they may be the same force.

These correlations however do not account for the difference in propagation through mediums initially, as it is well known that sound cannot travel through a vacuum.

By expanding the comparison chart to encompass other fields, we find that many still share the same calculations. Each is symbolically represented as a separate field with the first variable denoting this, as they exist as separate forces in nature.

Due to its property of requiring a medium and its inability to propagate through space sound is often held at arm's length separately from these forces, but is still worth noting for its wave properties and calculations to better illustrate the behavior of the larger forces in our minds and imagination.

\section*{Shared Function and Wave Type of Various Field Calculations}

\begin{table}
    \centering
    \begin{tabular}{ccc}
         \textbf{Force (Amplitude)} & \textbf{Unique Variable} & \textbf{Shared $\propto$} \\
Gravity & $h$ & $\frac{1}{r}$ \\
Gravitational Potential & $\phi$ & $\frac{1}{r}$ \\
Sound & $A$ & $\frac{1}{r}$ \\
Electric Fields & $E$ & $\frac{1}{r}$ \\
Magnetic Fields & $B$ & $\frac{1}{r}$
    \end{tabular}
    \caption{Caption}
    \label{tab:Table 1}
\end{table}


Also of note, we find a correlation pertaining to propagation in the way we measure the speed of propagation for Gravity, Electricity, and Magnetism, all propagating at the speed of light. And although it is not as accurate, we can apply the same speed measurement to sound as well, denoting a difference in scale. As such we can also include sound on the comparison chart for propagation speeds to illustrate that there is commonality among wave forms from all walks of different forces.

The previously noted comparisons raise an interesting implication about much of the measurements and calculations. While our interpretations have placed these areas of study into separate fields, on a numerical, and mathematical base level they calculate and in most cases share properties that are all calculated the same, beyond the simple property of being a wave.

Given the success of applying Hooke’s law to gravity we find that it can form a sound framework in order to unify a way of calculating the forces of waves and the effect it has on matter through a similar expression:

\[
F = - c x
\]

where:

\begin{itemize}
    \item $F$ is the force of the fundamental,
    \item For:
    \begin{itemize}
        \item Light: $F = $ Radiation Pressure,
        \item Gravity: $F = $ Gravitational Force,
        \item Magnetism: $F = $ Lorentz Force,
        \item Heat: $F = $ Energy Transfer,
    \end{itemize}
    \item $c$ is the constant,
    \item For:
    \begin{itemize}
        \item Light: $c = $ Light Speed,
        \item Gravity: $c = $ Gravitational Constant,
        \item Magnetism: $c = $ Field Strength,
        \item Heat: $c = $ Material Density,
    \end{itemize}
    \item $x$ is the effect of the force on matter,
    \item For:
    \begin{itemize}
        \item Light: $x = (I \cdot p)$ where $I=$ Intensity, $p=$ Photon Momentum,
        \item Gravity: $x = \frac{m_1 m_2}{r^2}$,
        \item Magnetic: $x = (I \cdot L)$ where $I=$ Current and $L=$ Length,
        \item Heat: $x = $ Expansion,
    \end{itemize}
\end{itemize}

While we do observe differences in the way they affect materials and react to outside forces, the correlations in calculating do suggest that beyond the wave properties of each there is an underlying fundamental or mechanism at work.

While gravity shares the properties of propagation and speed with the other transverse waves, under observation and even calculation it’s clear that gravity exhibits influence and effects that appears more consistent with the behavior of longitudinal waves, and we find a great many more correlations in calculating it with sound.

Therefore, the author proposes that a longitudinal wave component may be present within the force of gravity.

Under that consideration, a possibility arises. One that can propagate even through the vacuum of space, exerting force on the space time curvature, and moving matter:

\section*{Hypothesis}

\begin{quote}
Over time gravitational waves excite the gravitational potential producing a longitudinal pressure wave that travels at the speed of light propagating over oscillating electric and magnetic fields.
\end{quote}

\section*{Methods}

Using the already available frameworks, we can calculate a combination of equations that work together to produce this effect.

For a plane gravitational wave propagating along the $x$ axis over $t$ time, it can be expressed as:

\[
h(t, x) = h_0 \sin(k x - \omega t)
\]

where:

\begin{itemize}
    \item $h(t,x)$ is the strain of the wave at position $x$ and time $t$,
    \item $h_0$ is the amplitude of the waveform,
    \item $k$ is the wave number (wavenumber),
    \item $\omega$ is the angular frequency.
\end{itemize}

Now considering a disturbance in the gravitational potential due to the gravitational waves, a time dependent component can be added to its normally static state:

\[
\Phi(t, x) = \Phi_0 + \partial \Phi(t, x)
\]

where:

\begin{itemize}
    \item $\Phi_0$ is the gravitational potential,
    \item $\partial \Phi(t,x)$ represents the perturbation of the gravitational potential due to the waves.
\end{itemize}

\section*{Resulting Longitudinal Pressure Wave Expression}

\[
g(t, x) = \rho \frac{\partial \Phi}{\partial t}
\]

where:

\begin{itemize}
    \item $g(t,x)$ is the pressure wave function at time $t$ and position $x$,
    \item $\rho$ is a proportionality constant or density,
    \item $\frac{\partial \Phi}{\partial t}$ is the time derivative of the disturbed gravitational potential.
\end{itemize}

Combining the gravitational waves and the impact on the gravitational potential the expression becomes:

\[
g(t,x) = \rho \frac{\partial}{\partial t} \left( \Phi_0 + \partial \Phi(t,x) \right)
\]

This expression shows how the time varying component of the gravitational potential, influenced by the gravitational waves, can result in the formation of a longitudinal pressure wave traveling in the direction of the gravitational wave propagation.

\section*{Results and Discussion}

Using the calculations and the capabilities of this platform, we can produce interactive visuals that provide a clear picture.

The excited Gravitational Potential surface, along with electric field oscillations (red), and magnetic field oscillations (blue) provide a wonderful way to interact with time to see how the longitudinal wave is given a way to propagate through a vacuum over the fields.

Here the gravitational waves and excited gravitational potential are shown as streaming arrows, and the resulting gravitational potential pressure wave is represented with varying opacities.

\section*{Conclusion}

The potential gravitational pressure waves (PGPW) hold (for lack of a better term) much potential to help explain the mechanisms and the outer workings of gravity.

In the spirit of science the author looks forward to seeing what other researchers may find in their exploration of this potential and hopes that it may not only improve our cosmological models, but also humanity's understanding as a whole.

Or perhaps not as it may be, for there are 999 ways to incorrectly make a light bulb, and every mistake is another step forward in finding a solution.

And that is where the journey of discovery begins.

\end{document}

